\section{Pătrate Frecvente}
\textit{Dificultate: 3/5} \\
\textit{Propunător: Tudor Condorache} \\
\textit{Developer: Tudor Condorache, Denisa Natalia Pogar} \\
\textit{First solve: qfasfqwdfadafqwafvqvqvf după 00:31:16}  \\

\begin{remark}
    Dacă o secvență include $k$ materiale diferite, atunci lungimea sa trebuie să fie de:
    \[
        1^{2} + 2^{2} + \ldots + k^{2} = \frac{k(k+1)(2k + 1)}{6} \leq 500\ 000
    \]
    
\end{remark}

Avem deci $k_{max} = 114$ lungimi diferite ce pot fi testate, pe care le putem precalcula la început. Vom lua o secvență dintre acestea de lungime aribtrară $l$ și vom vedea, pentru fiecare \textit{sliding window} de lungime $l$, dacă aceasta formează o secvență mulțumitoare.

Pentru a vedea dacă o secvență este mulțumitoare, vom ține un vector de frecvență $f$, unde $f_{i}$ va reține de câte ori apare materialul $i$ în secvența noastră.

Acum, ca o secvență de lungime $l = 1^{2} + 2^{2} + \ldots + k^{2}$ să fie mulțumitoare, am dori ca în vectorul $f$ să avem
\begin{align*}
    f_{item_{1}}   &= 1 \\
    f_{item_{2}} &= 4 \\
    f_{item_{3}} &= 9 \\
     & \vdots \\
    f_{item_{k}} &= k^{2}
\end{align*}

Este totuși dificil să verificăm existența acestor valori în timp constant. De aceea, vom mai ține un vector de frecvență $f'$, unde $f'_{i} = 1 \iff \exists j \; \text{a.î}\; f_{j} = i^{2}$. Acum, condiția noastră se transformă în:
\begin{align*}
    f'_{j_{1}}   &= 1 \\
    f'_{j_{2}} &= 1 \\
    f'_{j_{3}} &= 1 \\
    & \vdots \\
    f'_{j_{k}} &= 1
\end{align*}

Acest lucru este mult mai ușor de verificat. Atunci când tranziționăm cu secvența de la $[l, r]$ la $[l + 1, r +1]$, trebuie să modificăm doar felul în care $A[r + 1]$ și $A[l]$ contribuie la cei doi vectori de frecvență.

\medskip

Complexitate finală: $\mathcal{O}(N \cdot k_{max})$.
